\documentclass[10pt]{article}

\usepackage[utf8]{inputenc}

\usepackage[T1]{fontenc}

\usepackage{amsmath}

\usepackage{amsfonts}

\usepackage{amssymb}

\usepackage[version=4]{mhchem}

\usepackage{stmaryrd}

\usepackage{mathrsfs}

\usepackage{bbold}



\begin{document}

щая собой множество преобразований $\{a, \Lambda\}$ с $\Lambda \in L_{+}^{\uparrow}$, содержащая единичное преобразование. В дальнейшем речь будет идти именно об этой группе.



Группа $\mathscr{P}_{+}^{\uparrow}-10$-параметрическая; к шести генераторам $M_{\mu \nu}$ группы Лоренца добавляются четыре генератора $P_{\mu}$ трансляций. Алгебра Ли П. г. определяется перестановочными соотношениями для генераторов:



\[

\begin{gathered}

{\left[M_{\mu v}, M_{\rho \sigma}\right]=i\left(g_{\mu \rho} M_{v \sigma}-g_{\mu \sigma} M_{v \rho}-g_{v \rho} M_{\mu \sigma}+g_{v \delta} M_{\mu \rho}\right),} \\

{\left[P_{\mu}, P_{v}\right]=0,\left[M_{\mu v}, P_{\rho}\right]=i\left(g_{\mu \rho} P_{v}-g_{v \rho} P_{\mu}\right),}

\end{gathered}

\]



где $g_{\mu \nu}-$ метрич. тензор. Десять генераторов П. г. являются основными динамич. величинами в релятивистской механике. Величину $P_{\mu}$ называют вектором энер ги и - имп у л ь с а или 4-и м п у л ьс о м; 3-вектор $\boldsymbol{M}=\left(M_{23}, M_{31}, M_{12}\right)$ есть угловой момент. В квантовой теории поля для любого оператора $A(x)$



\[

\left[A(x), P_{\mu}\right]=i \partial A(x) / \partial x^{\mu} .

\]



В частности, эволюция во времени определяется оператором $P_{0}$, или гамильтонианом системы.



Для П. г. имеются два Казимира оператора, коммутирующих со всеми ее генераторами и, следовательно, релятивистски инвариантных. Это $P^{2}=P_{\mu} P^{\mu}$ и $W=-w^{2}$, где псевдовектор $w^{\mu}=(1 / 2) \varepsilon^{\mu \nu \lambda \sigma} M_{v \lambda} P_{\sigma}$, а $\varepsilon^{\mu \nu \lambda \sigma}-$ полностью антисимметричный тензор.



При $P^{2} \geqslant 0$ имеется еще одна дискретная инвариантная характеристика - знак энергии: $\varepsilon=P_{0} /\left|P_{0}\right|$ с собственными значениями \textbackslash pm 1 .



Как и в случае группы Лоренца, представления П. г. строят с помощью односвязной группы $\mathscr{P}_{0}-$ универсальной накрывающей для группы $\mathscr{P}_{+}^{\uparrow}$. Для квантовой теории поля важны унитарные неприводимые представления $\mathscr{P}_{+}^{\uparrow}$. Сoгласно требованию релятивистской инвариантности, векторам состояния отвечают так наз. проективные представления, задаваемые с точностью до фазового множителя. Имеет место те орема Вигнера - Баргмана, утверждающая, что любое проективное представление группы $\mathscr{P}_{+}^{\uparrow}$ порождается обычным однозначным унитарным представлением группы $\mathscr{P}_{0}$.



Изучение важных для физики унитарных представлений группы $\mathscr{P}_{0}$ сводится к классификации ее неприводимых унитарных представлений, так как хотя $\mathscr{P}_{0}$ и некомпактна, любое ее унитарное представление может быть разложено в прямую сумму (или интеграл) неприводимых представлений.



Группа $\mathscr{P}_{0}$ локально изоморфна группе $\mathscr{P}_{+}^{\uparrow}$ и имеет те же генераторы и те же операторы Казимира, что и $\mathscr{P}_{+}^{\uparrow}$. В зависимости от значений оператора $P^{2}$ представления группы $\mathscr{P}_{0}$ могут быть разделены на следующие классы:



\begin{enumerate}

  \item $P^{2}=m^{2}>0$.

\end{enumerate}



la) $\varepsilon=1$ (то есть $P_{0}>0$ ). Соответствующие представления описывают трансформационные свойства реальных частиц с массой покоя $m$.



1б) $\varepsilon=-1$ (то есть $P_{0}<0$ ). Эти представления комплексно сопряжены с представлениями класса 1a).



\begin{enumerate}

  \setcounter{enumi}{1}

  \item $P^{2}=0, P \not 0$.

\end{enumerate}



2a) $\varepsilon=1\left(P_{0}>0\right)$. Соответствующие представления описывают частицы с нулевой массой покоя (нейтрино и фоТон).



2б) $\varepsilon=-1\left(P_{0}<0\right)$. Представления этого класса комплексно сопряжены с представлениями класса 2a).



\begin{enumerate}

  \setcounter{enumi}{2}

  \item $P^{2}=-m^{2}<0$ (то есть вектор $\boldsymbol{P}$ пространственноподобен). Согласно основным принципам релятивистской механики, частицы с таким импульсом не могут реально суще- ствовать. Однако представления класса 3) также встречаются в квантовой теории поля, напр. при описании трансформационных свойств взаимодействующих полей.



  \item $\boldsymbol{P}=0$. Все состояния с таким $P$ трансляционно инвариантны. Все унитарные представления этого класса, кроме единичного, бесконечномерны. Единичное представление соответствует вакууму, инвариантному относительно всех преобразований из П. г.



\end{enumerate}



Физич. смысл инварианта $w^{2}$ выявляется просто при $m^{2}>0, P_{0}>0$. В этом случае величина $-w^{2} / m^{2}$ равна квадрату углового момента $\boldsymbol{M}^{2}$ в состоянии покоя, то есть квадрату спина.



Таким образом, неприводимое унитарное представление П. г. характеризуется значениями массы $m$, спина $S$ и знака энергии (при $m^{2}>0$ ).



Jum.: [1] Боголюбов Н.Н., Логунов А.А., Тодоро в И.Т., Основы аксиоматического подхода в квантовой теории поля, М., 1969; [2] Н о в о ж и л о в Ю. В., Введение в теорию элементарных частиц, М., 1972; [3] М ишель ь., Ша аф М., Симметрия в квантовой физике, пер. с англ., М., 1974; [4] Бару т А., Рон ч к а Р., Теория представлений групп и ее приложения, пер. с англ., т. 1-2, М., 1980; [5] Элл и от Дж., Доб е р П., Симметрия в физике, пер. с англ., т. 1-2, М., 1983. ПУАНКАРЕ ЗАДАЧА - см. Краевая задача теории аналитических функций.



ПУАНКАРЕ ИНДЕКС о о обо й точк и - топологический инвариант изолированной особой точки $O$ непрерывного векторного поля $\boldsymbol{V}$ на $n$-мерном многообразии $M$, равный, по определению, степени отображения $f: S \rightarrow S$, заданного формулой



\[

f(x)=\rho V(x) /|V(x)|,

\]



где $x$ - система локальных координат на $M$ с центром в $O$, a $S-$ сфера $|x|=\rho$ достаточно малого радиуса $\rho$. П. и. особой точки $O$ поля $V$ на плоскости $(n=2)$ совпадает с числом оборотов вектора $V$ при обходе вокруг $O$ по окружности достаточно малого радиуса. П. и. невырожденной особой точки $O$ гладкого векторного поля $V$ равен $\operatorname{sign}\left(\lambda_{1} \lambda_{2} \ldots \lambda_{n}\right)$, где $\lambda_{1}, \lambda_{2}, \ldots, \lambda_{n}$ - собственные числа линеаризации поля $\boldsymbol{V}$ в $\boldsymbol{O}$. В частности, при $n=2$ П. и. узла, фокуса, центра и центро-фокуса равен 1, а П. и. седла равен -1 .



Понятие индекса особой точки, введенное А. Пуанкаре в 1886 (см. [1]) (только для невырожденных особых точек), восходит, так же как и более общее понятие степени отображения, к так наз. характеристике Кронекера системы функций (см. [2]), поэтому П. и. часто называют и ндек сом Кронекера - Пуанкаре (см. [6]).



Если все особые точки векторного поля $\boldsymbol{V}$ на компактном многообразии без края изолированы, то сумма их П. и. не зависит от $V$ и равна эйлеровой характеристике многообразия (теорема Пуанкаре, или теорема Пуанкаре - Xопфа).



Лит.: [1] П у а к а р е А., О кривых, определяемых дифференциальными уравнениями, пер. с франц., М.-Л., 1947; [2] Kronecker L., "Monatsber. Berl. Akad.», 1869, S. 159-93, 688-98; [3] H opf H., «Math. Ann.», 1927, Bd 96, S. 225-50; [4] Apн о л ь Д В. И., Обыкновенные дифференциальные уравнения, 3 изд., М., 1984; [5] Дубровин Б.А., Новиков С. П., Фоменко А. Т., Современная геометрия, 2 изд., М., 1986; [6] Аносов Д. В. [и др.], в кн.: Итоги науки и техники. Современные проблемы математики. Фундаментальные направления, т. 1, Динамические системы, М., 1985, с. 151-242; [7] Х и рш М., Дифференциальная топология, пер. с англ., М., $1979 . \quad$ М. Б. Севрюк. ПУАНКАРЕ ИНТЕТРАЛЬНЫЙ ИНВАРИАНТ - ДИФФеренциальная 1-форма $p d q$ в стандартном симплектическом пространстве $\left(\mathbb{R}^{2 n}, d p \wedge d q\right)$. Фазовый поток гамильтоновой системы сохраняет интеграл от П. и. и. по замкнутым кривым:



\[

\int_{g \gamma} p d q=\int_{\gamma} p d q

\]



М. Э. Казарян.





\end{document}